\documentclass[11pt]{article}

% --- packages -------------------------------------------------------------
\usepackage[a4paper,margin=1in]{geometry}
\usepackage[T1]{fontenc}
\usepackage{newtxtext,newtxmath}   % Times-like text + math; drop if unavailable
\usepackage{amsmath}   % newtxmath supplies AMS symbols; loading amssymb here clashes (\Bbbk)
\usepackage{authblk}               % clean multi-affiliation author block
\usepackage[hidelinks]{hyperref}
\usepackage{siunitx}

% --- author-block spacing tweaks -----------------------------------------
\renewcommand\Authfont{\bfseries}
\setlength{\affilsep}{0.4em}

% --- title ----------------------------------------------------------------
\title{\bfseries Harder, Better, Faster, Stronger Neutronics for Fusion
Reactors: Stellarator and Spin-Polarized Neutronics Sources in OpenMC and
Automated Shielding Optimization}

% --- authors and affiliations --------------------------------------------
\author[1,2]{Thaddaeus Kiker\thanks{Corresponding author:
\href{mailto:tjk2147@columbia.edu}{tjk2147@columbia.edu}}}
\author[3]{Ethan Peterson}

\affil[1]{Department of Astronomy, Columbia University,
550 West 120th Street, New York, NY 10027, USA}
\affil[2]{Columbia Fusion Research Center, Columbia University,
New York, NY 10027, USA}
\affil[3]{Plasma Science and Fusion Center, Massachusetts Institute of
Technology, 175 Albany Street, Cambridge, MA 02139, USA}

\date{\today}

% --- keywords macro -------------------------------------------------------
\providecommand{\keywords}[1]{\vspace{0.6em}\noindent\small\textbf{Keywords:} #1}

\begin{document}
\maketitle

\begin{abstract}
\noindent
Superconducting magnets are the most radiation-sensitive and costly
components of a fusion power plant, and protecting them from the intense
\SI{14}{\MeV} neutron flux requires shielding tuned inside Monte Carlo
transport simulations that can take hours per design. This work confronts a
\emph{harder} version of that problem than standard axisymmetric tools allow.
We add physics-informed fusion neutron sources to the open-source OpenMC code:
a native VMEC-fluxmap \texttt{StellaratorSource} that captures the true,
twisting three-dimensional geometry of stellarator and tokamak plasmas, and a
Schwartz-kernel birth-direction sampler that captures the anisotropic,
spin-steered emission of spin-polarized fuel. These sources give a
\emph{better}, higher-fidelity account of where neutrons are born and aimed:
the polarized kernel reproduces the analytic angular limits of Schwartz (2025)
and, unlike precomputed-point approaches, evaluates the local magnetic-field
direction natively and composes with any device geometry. Building on them, we
develop a \emph{faster}, direct-forward automated shield-optimization loop that
resolves the per-magnet neutron load to a few percent in minutes---removing the
variance-reduction machinery previously assumed necessary---and drives an
attribution-guided redesign of the shield. The outcome is \emph{stronger}
magnet protection from the same footprint: concentrating shielding where
neutron attribution directs it cut the peak coil fast-neutron load roughly
fourfold ($4.2\times$) versus a uniform baseline at fixed envelope. An exact
volume-calculation pipeline further reveals that identical nominal shield
thickness varies $2.4\times$ in mass by location, and that the inboard coils
sitting closest to the plasma---which see the highest neutron load and most
need shielding---are precisely where the least physical room exists to add it.
The controlling result is quantitative and non-obvious: because equal shielding
thickness costs up to $2.4\times$ different mass and the most heavily loaded
coils have the least room to grow, stellarator magnet protection is bounded by
geometric \emph{feasibility} rather than shielding mass---set by where a design
\emph{can} add shielding, not merely where it should.
\end{abstract}

\keywords{fusion neutronics; stellarator; spin-polarized fusion; OpenMC;
Monte Carlo neutron transport; superconducting-magnet shielding; DAGMC;
shield optimization; VMEC equilibria; tritium breeding}

\end{document}
